\documentclass{article}
\usepackage[left=2cm, right=2cm, top=2cm, bottom=2cm]{geometry}
\usepackage{graphicx}
\usepackage{amsmath}
\usepackage{amssymb}
\usepackage{amsthm}
\usepackage{fancyhdr}
\usepackage{verbatim}
\usepackage{listings}
\usepackage{xcolor}
\usepackage{pdfpages}
\usepackage{cancel}
\usepackage{physics}
\usepackage{esint}
\usepackage{braket}

\lstset{
    language=C++,
    basicstyle=\ttfamily\footnotesize,
    keywordstyle=\color{blue},
    commentstyle=\color{green},
    stringstyle=\color{red},
    numbers=left,
    numberstyle=\tiny\color{gray},
    frame=single,
    breaklines=true,
    captionpos=b,
}


\title{MTH 446: Lecture \# 17}
\author{Cliff Sun}

\newtheorem{theorem}{Theorem}[section]
\newtheorem{lemma}[theorem]{Lemma}
\newtheorem{definition}[theorem]{Definition}
\newtheorem{conjecture}[theorem]{Conjecture}
\newtheorem{proposition}[theorem]{Proposition}
\newtheorem{corollary}[theorem]{Corollary}
\newtheorem{one minute paper}[theorem]{One Minute Paper}

\pagestyle{fancy}
\lhead{\textbf{\thepage}\ \ \nouppercase{\rightmark}}
\chead{MTH 446: Lecture \# 17}
\rhead{Cliff Sun}

\begin{document}

\maketitle

\section*{Lecture Span}
\begin{enumerate}
    \item Complex logs
\end{enumerate}

\section*{Complex logs}

Consider the analytic single-valued branch of $\log z$ as 

\begin{equation}
    \log z = \ln r + i\theta, \quad \frac{3}{4}\pi < \theta < \frac{11}{4}\pi
\end{equation}

Recall that 

\begin{equation}
    (1-i)^{4i} = \exp(4i\ln(1-i))
\end{equation}
Then find 

\begin{equation}
    \left( \theta \in \arg(1-i) \right) \land (3\pi/4 < \theta < 11\pi/4)
\end{equation}

You have to choose the $n$ of which $3\pi/4 < \theta < 11\pi/4$. 

\section*{Analyticity of single valued of $z^c$}

We can prove that $z^c$ is analytic. Moreover, 

\begin{equation}
    \frac{d}{dz}(z^c) = \frac{c}{z}\exp(c\log(z)) = cz^{c-1}
\end{equation}

Moreover, we can define the following:

\begin{equation}
    c^z = \exp(z\log c)
\end{equation}

\section*{Trigonometric Functions}

\begin{definition}
    For every $z \in \mathbb{C}$, define 
    \begin{equation}
        \cos z = \frac{e^{iz} + e^{-iz}}{2}, \quad \sin z = \frac{e^{iz} - e^{-iz}}{2i}
    \end{equation}
    It can be proved that $\cos z$ and $\sin z$ are entire functions. Thus,  $\cos, \sin \in H(\mathbb{C})$. Such functions follow the standard trigonometric differentiation rules. 
\end{definition}

You can also show 

\begin{enumerate}
    \item $\sin(z_1 + z_2) = \sin z_1 \cos z_2 + \sin z_2 \cos z_1$
    \item $\sin(z_1 - z_2) = \sin z_1 \cos z_2 - \sin z_2 \cos z_1$
    \item $\cos(z_1 + z_2) = \cos z_1 \cos z_2 - \sin z_1 \sin z_2$
    \item $\cos(z_1 - z_2) = \cos z_1 \cos z_2 + \sin z_1 \sin z_2$
\end{enumerate}

Prove that $f(z) = (1 - \cos 2z )/2 = \sin^2 z = g(z)$.

\begin{proof}
    Since $f(z) = g(z)$ for $z \in \mathbb{R}$, by the uniqueness property, it follows that $f(z) = g(z)$ for all $z \in \mathbb{C}$. 
\end{proof}

You can also prove all the trig identities by the uniqueness property. We can also calculate 

\begin{equation}
    \sin(i) = \frac{e - \exp(-1)}{2}i
\end{equation}

Note that $|\sin(i)| > 1$. 

Also, define 

\begin{equation}
    \cosh(x) = \frac{\exp(x) + \exp(-x)}{2}, \quad \sinh(x) = \frac{\exp(x) - \exp(-x)}{2}
\end{equation}

Define 

\begin{equation}
    \cosh(z) = \dots, \quad \sinh(z) = \dots
\end{equation}

Note 

\begin{enumerate}
    \item $\sin(iy) = i\sinh(y)$
    \item $\cos(iy) = \cosh(y)$
\end{enumerate}

Calculate 

\begin{align}
    \cos(z) &= \cos(x)\cos(iy) - \sin(x)\sin(iy) \\
    &= \cos x\cosh y - i \sin x \sinh y
\end{align}

\begin{align}
    \cos^2 z &= \cos^2 x \cosh^2 y  + \sin^2 y \sin^2 x \\
    &= \cos^2x (1 + \sinh^2 y) + \sinh^2 y \sin^2x \\
    &= \cos^2 x + \sinh^2 y
\end{align}

Therefore, 

\begin{align}
    |\cos z| &= \sqrt{\cos^2 x + \sinh^2 y} \\
    |\sin z| &= \sqrt{\sin^2 x + \sinh^2 y}
\end{align}

Moreover, 

\begin{equation}
    \frac{d}{dz}\cosh z = \sinh z
\end{equation}



\end{document}